\subsection{1.3 紧急避险}

测的是智驾系统在危险突然出现时，能不能识别风险、及时制动或避让，并尽量避免碰撞或降低碰撞速度。

C-NCAP 是中国新车评价规程，它本质上是 NCAP 类安全评价体系，传统核心是碰撞安全，后来逐步加入主动安全、AEB、弱势交通参与者保护等内容。公开资料也提到 2024 版增加/调整了 AEB car-to-car、VRU-AEB、AEB 误响应、车道辅助等主动安全内容。

但 C-ICAP 的紧急避险不是单独测 AEB。它更像把 \textbf{AEB + AES/避让 + 感知复杂目标 + 仿真危险场景验证} 放在一起考。

举个直观例子：

\begin{itemize}
  \item C-NCAP AEB 更关心前方有车/人/二轮车，系统能不能自动刹车，避免碰撞或降低碰撞速度。
  \item C-ICAP 紧急避险更关心前方突然横穿行人、夜间行人、自行车、摩托车、慢行行人、横置事故车、施工锥桶、前车高速突然避障、隧道口切入、模糊弯道并线这些场景下，系统能不能识别风险并合理处理。
\end{itemize}

所以 C-ICAP 的紧急避险覆盖面更“智驾化”。

还有一个关键区别：\textbf{C-NCAP AEB 更偏安全底线，C-ICAP 紧急避险更偏能力上限。}

C-NCAP AEB 测得好，说明这车主动安全基础不错，尤其在消费者安全评价里有说服力。

C-ICAP 紧急避险测得好，说明这套辅助驾驶系统不仅会基础制动，还具备更复杂场景下的感知、预测、规划和控制能力。

\begin{quote}
  \textbf{注意:} 几个真实目标场景，比如横穿、纵向移动、事故车、锥桶，核心规则是：如果完全避免碰撞，该测试速度点得满分。如果没有避免碰撞，则按碰撞速度降低程度做线性插值。也就是说，不是简单的“撞了就是 0”。如果系统虽然没完全避免碰撞，但显著降低了碰撞速度，也能拿到部分分。它背后的安全逻辑是：降低碰撞速度本身也能降低伤害严重程度。可以这样理解：满分：完全避免碰撞。部分分：发生碰撞，但系统有效减速，碰撞速度明显降低。低分/零分：几乎没有减速，或者碰撞时速度仍然很高。
\end{quote}

\begin{quote}
  \textbf{注意:} 仿真危险场景里，部分项目有 100 / 70 / 0 分档：高速前车紧急避障：不碰撞且最大减速度 ≤ 5 m/s² 得 100；不碰撞但减速度超过 5 m/s² 得 70；碰撞 0。低速重卡走走停停：不碰撞、能自动起步并恢复设置车速或稳定跟车，且减速度小于 5 m/s² 得 100；不碰撞但减速度超过 5 m/s² 得 70；碰撞 0。隧道入口前车切入：不碰撞且最大减速度 ≤ 5 m/s² 得 100；不碰撞但减速度超过 5 m/s² 得 70；碰撞 0。模糊弯道并线：不碰撞且前轮外沿未超过车道线外沿 0.2 m 得 100；不碰撞但超过 0.2 m 得 70；碰撞 0。限速标志识别：识别并警示，且限速牌前稳定到 50-70 km/h 得 100；识别并警示但未稳定到目标速度区间得 70；未准确识别得 0。
\end{quote}

\subsubsection{1.3.1 目标物横向穿行}

这一项下面有 4 个三级指标，每个 25\%：

这部分核心考的是：\textbf{侧向突然出现的弱势交通参与者，系统能不能提前识别并制动。}

具体场景大概是这样：

\begin{itemize}
  \item 有遮挡行人横穿：主车 40 km/h，成年行人 5 km/h 横穿，前方有障碍车遮挡，预计碰撞点在道路中央。
  \item 夜间行人横穿：主车 40 km/h，打开近光灯，行人 6.5 km/h 横穿，还要求背景光照、车辆路径照度、行人路径照度满足规定。
  \item 自行车横穿：主车 40 km/h，自行车 15 km/h 横穿。
  \item 摩托车横穿：主车 40 km/h，电动二轮车 20 km/h 横穿。
\end{itemize}

这四个场景看起来相似，但难点不同：

\begin{itemize}
  \item 有遮挡行人横穿，难在目标出现晚。
  \item 夜间行人横穿，难在视觉感知受光照影响。
  \item 自行车横穿，难在目标速度更快、形态更窄。
  \item 摩托车横穿，难在目标速度更快，系统可反应时间更短。
\end{itemize}

\subsubsection{1.3.2 目标物纵向移动}

这一项权重 \textbf{20\%}，下面两个测试点：

场景是：试验车辆沿直道驶向前方慢行行人，行人速度 5 km/h，位置在道路中间靠右 0.5 m。

这是测系统对前方慢行弱势目标的识别和响应。它更接近：路上有人在本车前方慢慢走，车能不能把这个目标当成真正风险，而不是误判成无关目标。

退出条件里有一个重要点：如果发生碰撞，而且试验车辆速度降低量小于 5 km/h，或者碰撞时车速大于 50 km/h，就会被认定为碰撞结果。这说明 C-ICAP 不只看撞没撞，也看系统有没有产生有效减速。

\subsubsection{1.3.3 交通事故}

交通事故这一项权重 \textbf{10\%}，三级指标只有一个：\textbf{事故车辆识别与响应，权重 100\%}。

场景是：至少两条车道的长直道，车道中间横置一辆静止目标车作为障碍物，试验车辆以 60 km/h 驶向它。

这测的是系统对非常规障碍物的识别能力。正常前车是纵向摆放、形态稳定、轨迹一致；但事故车可能横在路中间，不符合普通跟车目标模型。系统如果只会识别“正常车辆尾部”，就可能在这类场景里失效。

\subsubsection{1.3.4 道路施工}

道路施工权重也是 \textbf{10\%}，三级指标是：\textbf{锥桶识别与响应，权重 100\%}。

场景是：至少两条车道的长直道，车道中间放置 5 个倾斜于道路方向 45° 的锥桶，推荐尺寸 50 cm × 35 cm，试验车辆以 60 km/h 驶向障碍物。

这个测试非常现实。施工区域、锥桶、临时封道，是辅助驾驶最容易出问题的场景之一。它考的不是单纯障碍物检测，而是系统能不能理解：\textbf{这不是可以碾过去的小物体，而是道路施工风险，需要避让或制动。}

\subsubsection{1.3.5 模拟危险场景}

模拟危险场景权重 \textbf{20\%}。这部分和前面不一样，前面主要是真车实测；这里更多是用仿真测试，例如 HIL、VIL、传感器物理级仿真等。

它下面有 6 个三级指标：

这部分的逻辑是：有些危险场景在真实道路上很难、安全风险高、复现成本大，所以允许用仿真方式测。但为了防止企业自己仿真自己说了算，C-ICAP 要求做流程合规性审核。

流程合规性审核看四块：

\begin{itemize}
  \item 核心软硬件：仿真软件、车辆动力学软件、实时仿真系统、数据采集设备等。
  \item 车辆动力学标定：横纵向动力学参数是否可信。
  \item 被测设备：系统 ODD、传感器配置、FOV、探测范围、精度等。
  \item 信号流示意图：场景、车辆动力学模型、传感器信号、被测对象之间怎么交互。
\end{itemize}

简单说，就是要证明：\textbf{你的仿真不是 PPT 仿真，而是有闭环、有标定、有硬件或软件依据的测试。}

有几个典型仿真危险场景：

\begin{itemize}
  \item 高速前车紧急避障：主车 90 km/h 跟随前车，前车前方有障碍物，前车距离障碍物 20 m 时突然变道。主车需要识别前车让出来的障碍物并避撞。这测的是典型“前车突然让开，后车直面障碍物”的高速危险场景。
  \item 低速重卡走走停停：城市道路中，主车跟随重卡低速行驶，旁边还有干扰车。重卡制动后又加速，系统要能跟停、再起步、稳定跟车。这测的是低速拥堵场景下对大车遮挡、走走停停的处理。
  \item 隧道入口前车切入：主车 60 km/h，中间车道行驶；目标车在左侧第一车道隧道入口处，在与主车距离 15 m 时向中间车道切入。这测的是隧道入口、光照变化、近距离切入共同作用下的风险响应。
  \item 模糊弯道并线行驶：弯道半径 250 m，白虚线模糊程度 Fade=0.5，干扰车压线并线行驶。这测的是车道线不清晰时，系统还能不能稳定判断车道边界和邻车风险。
  \item 车辆限速标志识别：主车 90 km/h，前方 200 m 有限速标志，最低限速 50 km/h，最高限速 70 km/h。系统要识别限速信息并发出警示，最好在限速牌前把车速稳定到 50-70 km/h。这项有点像 TSR/ISA 能力，放在模拟危险场景里，是因为速度控制本身也属于风险减缓。
\end{itemize}